% =============================================================================
% Chapter 1 — Introduction
% Tara Torbati — Master's Thesis
% =============================================================================

\chapter{Introduction}
\label{chap:intro}

\section{Motivation and Context}
\label{sec:intro-motivation}

Agriculture consumes roughly 70\% of global freshwater
withdrawals~\cite{Sharma2025,Bayar2025}, and that share is concentrated
in regions where water is already scarce. Iran sits at the extreme end
of this spectrum: a combination of long-term aridification, decades of
supply-side infrastructure expansion, and a regulatory framework that
prices irrigation water at levels two to three orders of magnitude
below industrial and domestic tariffs has produced a hydrological
system in which the agricultural sector accounts for over 90\% of
withdrawals while being the user category least exposed to
volumetric-cost discipline~\cite{Mesgaran2018,Nouri2023,Behling2022}.
The consequences of this asymmetry are visible at multiple scales:
groundwater levels declining at meters per year across major aquifers,
rivers desiccating before reaching the Caspian and Persian Gulf
basins, land subsidence on the order of centimeters per year in the
central plateau, and a widening structural gap between water demand
and water supply that conventional supply-augmentation strategies have
been unable to close~\cite{PerezBlanco2020,Hajirad2025}.

These dynamics are particularly acute in Gilan Province on Iran's
Caspian coast, the country's principal rice-producing region. Hashemi
rice, the dominant cultivar grown in Gilan, is irrigated through
traditional flooded-paddy practice that consumes water at rates well
above the Penman--Monteith evaporative demand of the
crop~\cite{jalali2021,sadidi2021}. Subsidized agricultural water
pricing eliminates the economic feedback that would otherwise
discipline this consumption. The result is a regime in which
individual farmers face no marginal cost for over-application, while
the cumulative system-level consequences --- aquifer depletion, soil
salinization, and reduced downstream
flow~\cite{PourgholamAmiji2024,Behling2022} --- are externalized
across the basin and across time.

The technical-control framing of this problem treats irrigation as a
constrained finite-horizon optimal control problem: allocate a finite
seasonal water budget across a 90+ day cultivation season to maximize
crop yield while respecting per-day actuator capacities, soil-physical
state bounds, and a hard cumulative volumetric constraint. This
framing does not require policy reform as a precondition for
deployment: a controller that respects a season-cumulative budget can
be deployed under any tariff schedule, including the current
subsidized one, and the budget itself can be set externally by water
resource authorities, by aquifer-level monitoring, or by regional
allocation policy. The technical question is whether the controller
can deliver agronomically acceptable yields given a finite water
allocation, and how it should distribute that allocation in space
and time across a topographically heterogeneous field.

\section{Problem Statement}
\label{sec:intro-problem}

This thesis develops and compares constrained optimal-control and
reinforcement-learning approaches to irrigation on a 130-agent
topographically heterogeneous Hashemi rice field in Gilan Province,
Iran. The field is represented as a discrete grid of cells with
non-trivial elevation gradients (111~m of vertical relief over the
field extent), each cell maintaining its own coupled soil-moisture,
phenological, and biomass states, and exchanging water laterally
through topographically driven surface runoff. Each cell admits an
independent irrigation control input bounded by the physical actuator
capacity, with the joint control vector subject to a global
season-cumulative water budget representing the realistic deficit
allocations issued under regional water-management
schemes~\cite{Mesgaran2018,Nouri2023}.

Within this setting, the thesis investigates four specific questions:

\begin{enumerate}
\item \textbf{Cost-function design.} How should the multi-objective
cost function of a constrained MPC controller be structured and
weighted to balance terminal biomass, water consumption, drought
risk, and waterlogging avoidance, and how should those weights be
calibrated to real economic price signals?

\item \textbf{Pathology characterization.} What dynamic pathologies
arise in the controller-plant interaction at long prediction
horizons, in saturated-soil regimes, and under cheap water, and how
can these be resolved through cost-function design rather than
through ad-hoc heuristic correction?

\item \textbf{MPC versus RL comparison.} Under a fair, exact-equivalent
reward formulation, does a Soft Actor-Critic agent trained on the
ABM reproduce the policy quality of constrained MPC, and at what
trade-off between training cost and deployment latency?

\item \textbf{Robustness to scenario and budget.} How do both
controllers respond to climatic scenario variation (dry, moderate,
wet) and to deficit budget allocations (100\%, 85\%, 70\% of the
climatologically derived full requirement)?
\end{enumerate}

The answers to these questions are intended to inform two distinct
audiences: control-systems researchers studying constrained
distributed-parameter control, and water-management practitioners
considering technical deployment options under volumetric
deficit-irrigation policy.

\section{Contributions}
\label{sec:intro-contributions}

The principal contributions of this thesis are as follows.

\paragraph{A validated topographical agent-based virtual plant.}
A 130-agent ABM of a Gilan rice field is developed with explicit
two-layer hydrological decoupling between surface ponding and
root-zone moisture, a padded digital-elevation-model boundary
condition that resolves a degenerate sink behavior in prior
formulations~\cite{Lopezjimenez2024}, and crop-physiology stress
functions calibrated to Hashemi rice. The model is parameterized
from the FAO Irrigation and Drainage Paper No.~56
\cite{allen1998}, regional agronomic
literature~\cite{jalali2021,sadidi2021}, and the soil-hydraulic
database of Rawls et al.~\cite{rawls1982}. The plant constitutes a
fair, physics-grounded test bed on which both the predictive and the
learning-based controllers are evaluated.

\paragraph{A constrained MPC controller with economically anchored
cost weights.} A nonlinear MPC formulation is developed for the ABM,
solved via the IPOPT interior-point method within the CasADi symbolic
optimization framework. The cost function is calibrated to the
four-tier Iranian water-pricing schedule (subsidized agricultural,
domestic base, domestic high-consumption, and industrial), enabling
direct interpretation of the controller's behavior as a response to
real economic price signals rather than as a response to abstract
regularization weights.

\paragraph{A systematic weight-sensitivity analysis identifying and
resolving three controller pathologies.} A thirty-one-configuration
sensitivity sweep is conducted across one-at-a-time, joint factorial,
and cross-scenario validation phases. The sweep characterizes three
pathologies of the original cost-function formulation: silent
over-irrigation under cheap water, field-capacity overshoot at long
prediction horizons, and chronic waterlogging in wet-year scenarios.
The sweep further demonstrates that two of the originally proposed
penalty terms (surface ponding and field-capacity overshoot) are
dynamically redundant once one is sufficiently strong, motivating a
parsimonious five-term cost function. The recommended operating
point is identified and validated across all combinations of
scenario, budget, and prediction horizon.

\paragraph{A direct MPC--SAC comparison under exact reward
equivalence.} A Soft Actor-Critic agent is trained against a reward
function that is the exact mathematical negation of the MPC path
cost, enabling rigorous policy-quality comparison rather than the
more common algorithm-vs-algorithm benchmarking. The SAC architecture
employs Centralized Training with Decentralized Execution
(CTDE)~\cite{amato2024ctde}, with parameter sharing across the 130
spatial agents in the actor and centralized state-action evaluation
in a twin critic.

\paragraph{Comparative results across scenario, budget, and horizon.}
All four controllers (no-irrigation baseline, fixed-schedule
heuristic, constrained MPC, and SAC) are evaluated on a 9-cell
factorial grid spanning three climate scenarios (dry 2022, moderate
2020, wet 2024) and three budget allocations (100\%, 85\%, 70\% of
the climatological full demand). The results identify the operational
regimes in which each architecture excels, quantify the
yield-vs-water-vs-latency trade-off, and provide concrete guidance
on architecture selection as a function of deployment context.

\section{Scope and Limitations}
\label{sec:intro-scope}

The thesis operates within a deliberately bounded scope. Several
adjacent questions are out of scope and are noted here for clarity.

\textbf{Empirical field-level validation.} The agent-based plant is
parameterized from validated agronomic and pedological literature,
but empirical validation against direct field observations in Gilan
is not undertaken. The contribution is methodological: the comparison
of constrained controllers on a physics-grounded, parameter-validated
test bed. Empirical field deployment and post-deployment validation
are identified as natural directions for future work
(Chapter~\ref{chap:conclusion}).

\textbf{Sensor and actuator hardware.} The thesis treats per-cell
irrigation actuation as ideal --- requested volumes are applied
exactly, subject only to per-day capacity bounds. The practical
realization of cell-level variable-rate irrigation through commercial
hardware (modulated drip lines, modified center-pivot systems with
zone control, autonomous ground robots) is not addressed.

\textbf{Forecast uncertainty.} The principal experimental campaign
operates under perfect-forecast conditions to establish the upper
performance bound of the predictive architecture. A noisy-forecast
extension is described in Chapter~\ref{chap:controllers} and
evaluated in Chapter~\ref{chap:results}, but the full characterization
of forecast-uncertainty robustness across the joint scenario-budget-
horizon grid is left for future work.

\textbf{Multi-crop generalization.} The thesis is parameterized for
Hashemi rice. The agent-based framework supports other crops through
parameter substitution, but cross-crop validation is not undertaken.

\textbf{Economic and adoption analysis.} The water-pricing weights in
the MPC cost function are calibrated to real Iranian tariffs, but the
thesis does not undertake the broader cost-benefit analysis that
would be required to support a deployment recommendation.
Capital costs, sensor infrastructure, edge-computing hardware, and
adoption barriers are surveyed in the literature
(Chapter~\ref{chap:lit}) but not quantitatively analyzed.

\section{Thesis Structure}
\label{sec:intro-structure}

The remainder of this thesis is organized as follows.

\textbf{Chapter~\ref{chap:lit}} reviews the literature on irrigation
control, optimal and predictive control, reinforcement learning,
spatial flow modeling, and supporting infrastructure. The chapter
concludes with a synthesis identifying the specific research gap
addressed by the thesis.

\textbf{Chapter~\ref{chap:system}} presents the agent-based virtual
plant: the study site and digital elevation model, the directed
flow graph governing lateral water exchange, the climate data and
scenario selection, the soil and crop parameterization, the complete
state-space dynamics of the ABM, the action space and water-budget
constraint, and the unified software evaluation framework that
ensures fair comparison across controllers.

\textbf{Chapter~\ref{chap:controllers}} introduces the four controller
architectures: the no-irrigation rainfed baseline, the fixed-schedule
heuristic, the constrained Model Predictive Controller, and the Soft
Actor-Critic reinforcement learning agent. The MPC formulation is
developed in detail, including the multi-objective cost function with
economic anchoring, the symbolic control-oriented model with
smooth-approximation handling of non-differentiable dynamics, the
NLP construction and IPOPT solver configuration, the prediction-
horizon selection, and the warm-starting strategy. The chapter also
presents the systematic weight-sensitivity analysis that calibrates
the cost-function weights and resolves the three pathologies of the
initial formulation. The SAC architecture is developed with its
Markov decision process formulation, its CTDE actor-critic structure,
and its three-tier soft constraint-handling mechanism for the
seasonal water budget.

\textbf{Chapter~\ref{chap:results}} presents the comparative results
across the 9-cell factorial grid of climate scenario and budget
allocation, evaluating yield, water consumption, soil moisture
trajectories, and stress diagnostics for each of the four controllers
under perfect and noisy forecast conditions, and discussing the
identified operational trade-offs.

\textbf{Chapter~\ref{chap:implementation}} briefly addresses the
considerations relevant to physical implementation, including the
sensing modalities, communication protocols, and actuator hardware
that would be required to realize the controllers in deployment.

\textbf{Chapter~\ref{chap:conclusion}} summarizes the principal
findings, identifies the limitations of the present work, and
proposes directions for future research --- including empirical
field validation, robustness under forecast uncertainty,
generalization to additional crops and climatic regions, and the
integration of policy-driven economic feedback into the control
architecture.
